\chapter{Customization}

loclass verfolgt das Prinzip, sinnvolle Standards bereitzustellen, den Anwender jedoch nicht unnötig einzuschränken. Nahezu alle projektspezifischen Anpassungen erfolgen daher im Ordner \texttt{project}.

Während die Ordner \texttt{content} und \texttt{assets} den eigentlichen Dokumentinhalt aufnehmen, können im Ordner \texttt{project} individuelle Erweiterungen vorgenommen werden.

So können beispielsweise zusätzliche Pakete eingebunden, eigene Makros definiert oder neue Umgebungen erstellt werden. Diese stehen anschließend im gesamten Dokument zur Verfügung.

Vorhandene Befehle können bei Bedarf mit den üblichen \LaTeX{}-Mechanismen, beispielsweise mittels \texttt{\textbackslash renewcommand}, angepasst oder erweitert werden.

Der Ordner \texttt{core} sollte hingegen nur in Ausnahmefällen geändert werden. Er enthält die eigentliche Programmlogik von loclass. Anpassungen innerhalb dieses Verzeichnisses sind in der Regel nur dann erforderlich, wenn sich eine gewünschte Funktion nicht über die vorgesehenen Erweiterungsmöglichkeiten realisieren lässt.

Grundregel: Anpassungen erfolgen möglichst im Ordner project. Änderungen am core sollten die Ausnahme bleiben.
